\section{Model Validation}
\label{sec:validation}

Three experiments probe the RFSV pricer: (i) the implied volatility smile
against live SPY option chain data; (ii) a comparison with constant-volatility
baselines, and the roughness premium with and without variance matching; and (iii)
a parameter sensitivity sweep over $H$ and $\nu$.  Each uses the Python RFSV engine
(\texttt{data/rfsv\_model.py}), which implements the circulant-FFT path
generator with the same normalization as the C++ FFT pricer.  Every comparison
between configurations (different $H$, $\nu$ or $K$) uses common random numbers, so
differences are measured far more precisely than the absolute prices.

\subsection{Implied volatility smile}
\label{sec:iv}

\begin{figure}[t]
  \centering
  \includegraphics[width=\linewidth]{validate_iv.png}
  \caption{RFSV vs market implied volatility smile for SPY, for the two nearest
    weekly expirations in a snapshot taken on October 2, 2026 (spot \$771.09).
    Scatter: market IV computed by Black--Scholes inversion of mid-quote
    (\texttt{brentq}, $r = 0$, filtered to strikes in $[0.60S, 1.40S]$ with
    valid bid/ask spread); the $y$-axis is capped at $40\%$, and the 39 and 42
    deep-ITM market IVs above it are marked as off scale.
    Line: RFSV IV, same inversion applied to Monte Carlo European call prices,
    with common random numbers across strikes and the sample mean of $S_T$
    matched to the forward.
    The RFSV log-vol level $\mu_0 = \log\sigma_\mathrm{eff}$ is calibrated per expiration to
    match the ATM price (not IV): Brent's method on the price residual
    $p_\mathrm{RFSV} - p_\mathrm{mkt}(\sigma_\mathrm{ATM})$.
    The shaded region marks deep-ITM strikes where no simulated path finishes
    below $K$, so the Monte Carlo time value is zero and the IV is undefined.
    Controlled: $H = 0.10$, $\nu = 0.52$, $r = 0$, $M = 20{,}000$,
    $N = 63$ steps per year (at least 5 steps).
    Independent: strike $K$, expiration $T$.
    Dependent: implied volatility $\sigma_\mathrm{IV}$.}
  \label{fig:iv}
\end{figure}

Figure~\ref{fig:iv} shows the results.  For this comparison the base level is not
fixed at $\sigma_0 = 0.2$ but calibrated per expiration, $\mu_0 = \log\sigma_\mathrm{eff}$,
so that the RFSV ATM price matches the market ATM price.  This aligns the curves at
the money ($12.1\%$ and $13.1\%$ for the two expirations, with
$\sigma_\mathrm{eff} = 0.112$ and $0.119$), so the comparison concerns smile shape
(slope and curvature), not the vol level.

The model smile has clear curvature but no skew.  For the nearer expiration
($T = 0.019$) it has its minimum at the money and rises to $\approx 17.5\%$ at
$K/S = 0.91$ and $\approx 15\%$ at $K/S = 1.06$.  The SPY smile is strongly skewed
instead: above $40\%$ at $K/S = 0.91$, falling to a minimum of $\approx 10.5\%$ near
$K/S = 1.02$, and recovering to $\approx 16\%$ at $1.06$.  The later expiration
($T = 0.038$) has the same shape.  In the out-of-the-money call wing the two curves
roughly meet, so the vol-of-vol $\nu = 0.52$ produces about the right amount of
curvature there.  Below the money the model misses the skew entirely.  This is
expected: the engine draws the price shocks independently of the volatility driver,
so the spot-vol correlation is $\rho = 0$, and with $\rho = 0$ a stochastic volatility
model produces a smile that is symmetric in log-moneyness.  The steep short-dated
skew that rough volatility is known for \citep{gatheral2018volatility} requires
$\rho < 0$ (§\ref{sec:rho-future}).

\paragraph{A Monte Carlo pitfall: forward error masquerading as skew.}
An earlier run of this experiment (an October 1 snapshot with $\nu = 0.30$) used
$M = 3{,}000$ paths with common random numbers but no correction, and its RFSV curve rose steadily with strike, from
$13.0\%$ to $15.7\%$ across $0.98 \leq K/S \leq 1.05$.  That curve looks like an
upward skew, which a $\rho = 0$ model cannot produce.  The cause was sampling
error in the forward.  The sample mean of $S_T$ was $0.048\%$ below $S_0$, and
because every strike shares the same paths, that error acts like a shifted
forward and tilts the entire IV curve; it also biases the calibrated $\mu_0$.  A
400{,}000-path run of the same model gives an approximately symmetric smile with its minimum at
the money ($14.57\%$ at $K/S = 1.00$, against $14.79\%$ at $0.98$ and $14.70\%$
at $1.02$).  The current script rescales the simulated $S_T$ so that its sample
mean equals the forward exactly, a standard moment-matching correction
\citep{glasserman2004monte}, and uses $M = 20{,}000$ paths.  At $N \approx 5$
time steps per path the whole experiment runs in a few seconds.

\paragraph{Limitations}
I set $r = 0$ and ignore SPY's dividend yield.  For near-the-money options this
shifts IV levels slightly and leaves the shape comparison largely unaffected.
For deep in-the-money calls the effect is large: the carry in an ITM call price
is misread as time value, which inflates those market IVs (many above $40\%$),
so I do not interpret the low-strike wing.  SPY options are American-style; for
calls on a low-dividend ETF the early-exercise premium is small, and I treat
them as European.  The comparison is a single-day snapshot; market IV smiles
vary substantially day-to-day.  Brent's method stops at a tolerance of $0.01$ in
$\mu_0$, about $1\%$ in $\sigma_\mathrm{eff}$.

\subsection{Lévy benchmark and roughness premium}
\label{sec:levy}

\begin{figure}[t]
  \centering
  \includegraphics[width=\linewidth]{validate_asian.png}
  \caption{Asian option prices vs strike.
    \textbf{(a)} Price vs $K$: Lévy (1992) approximation and exact GBM by Monte
    Carlo at $\sigma = 0.2$ (constant vol), and RFSV at four $H$ values
    ($\nu = 0.52$, $\sigma_0 = 0.2$); each curve is the mean of 10 seeds of 10{,}000
    paths, with the same seeds for every curve.
    \textbf{(b)} Roughness premium $\hat{p}(H{=}0.10) - \hat{p}(H{=}0.50)$ vs $K$, at
    fixed $\nu = 0.52$ (red) and with $\nu$ at $H = 0.5$ chosen to match the expected
    integrated variance (blue); error bars are $\pm 1$ paired standard error across
    seeds.
    Controlled: $S_0 = 100$, $T = 1$, $\sigma_0 = 0.2$, $r = 0$, $N = 252$.
    Independent: $K$ (both panels); $H$ and $\nu$.
    Dependent: arithmetic Asian call price.}
  \label{fig:levy}
\end{figure}

Figure~\ref{fig:levy} shows the results.

\paragraph{Constant-volatility baselines: }
When $\nu \to 0$, $\sigma_t \to \sigma_0$, so the RFSV dynamics reduce to
constant-volatility GBM with $\sigma = 0.2$, which the Lévy approximation is designed
for \citep{levy1992pricing}.  At this volatility it is accurate: at the money it gives
$4.625$, against $4.606 \pm 0.027$ for exact discrete GBM by Monte Carlo
(100{,}000 paths).  All four RFSV curves lie above both baselines, ordered by $H$, as
stochastic volatility and Jensen's inequality predict.  At the money the RFSV prices
are $5.445$, $5.302$, $4.979$, and $4.834$ for $H = 0.05$, $0.10$, $0.30$, and $0.50$, so
rough stochastic volatility adds $0.70$ to the GBM price at $H = 0.10$.  (An earlier
version of this experiment used $\sigma = 1$, where the Lévy approximation overprices
exact GBM by about 6\%, enough to put the Lévy curve above the RFSV curves; it is
reliable only at realistic volatility.)

\paragraph{Roughness premium: }
At fixed $\nu = 0.52$, the premium $\hat{p}(H{=}0.10) - \hat{p}(H{=}0.50)$ is positive at
every strike, rising from $0.178$ at $K = 80$ to $0.468 \pm 0.008$ at the money and
falling to $0.319$ at $K = 120$.  But this compares more than roughness.  At fixed
$\nu$, $\text{Var}(\log\sigma_t) = \nu^2 t^{2H}$ is larger for small $H$ when $t < 1$, and
the expected integrated variance
\begin{equation}
  \label{eq:intvar}
  \mathbb{E}\!\left[\int_0^T \sigma_t^2\,dt\right]
  = \sigma_0^2 \int_0^T e^{2\nu^2 t^{2H}}\,dt
\end{equation}
is $18.6\%$ higher at $H = 0.10$ than at $H = 0.50$ ($1.574\,\sigma_0^2$ against
$1.327\,\sigma_0^2$).  To isolate roughness, I solve~\eqref{eq:intvar} for the $\nu$ at
$H = 0.5$ that gives the same integrated variance as $(H, \nu) = (0.10, 0.52)$, which
is $\nu = 0.651$, and price again with the same random numbers.  The variance-matched
premium is $0.330 \pm 0.009$ at the money, so about 70\% of the premium survives: most
of the price increase comes from path roughness itself, not from higher total
variance.  Both premia peak at the money, where the call payoff is most sensitive to
the shape of the volatility path.

\paragraph{Limitations: }
Matching expected integrated variance is one reasonable normalization, not the only
one; matching $\int_0^T \text{Var}(\log\sigma_t)\,dt$ instead gives a slightly different
$\nu$ (§\ref{sec:variancenorm}).  The standard errors are paired across 10 seeds, so
they reflect how consistently the premium appears under common random numbers; the
absolute prices carry the larger standard error of $\approx 0.03$.

\subsection{Parameter sensitivity: }
\label{sec:sensitivity}

\begin{figure}[t]
  \centering
  \includegraphics[width=\linewidth]{sensitivity_surface.png}
  \caption{ATM Asian call price vs model parameters.
    \textbf{(a)} Heatmap: price vs $(H, \nu)$ at $K = S_0 = 100$ (viridis
    colormap; low H = rough = lower on the $y$-axis because the axis is
    inverted so H increases upward).
    \textbf{(b)} Price vs $H$ for each $\nu$; curves decrease left-to-right
    and separate vertically with $\nu$.
    All cells share the same random numbers, so differences between cells are
    precise even though each absolute price has an MC standard error of
    $\approx 0.1$.
    Controlled: $S_0 = 100$, $K = 100$, $T = 1$, $\sigma_0 = 0.2$, $r = 0$, $N = 252$,
    $M = 10{,}000$.
    Independent: $H$, $\nu$.
    Dependent: ATM Asian call price.}
  \label{fig:sensitivity}
\end{figure}

\begin{figure}[t]
  \centering
  \includegraphics[width=\linewidth]{sensitivity_strike.png}
  \caption{Asian call price vs strike $K$ for varying $H$ at $\nu = 0.52$, with
    common random numbers across all curves.  Curves decrease monotonically with
    strike and are ordered by $H$ at every strike; the spread between $H = 0.05$
    and $H = 0.50$ peaks at the money.
    Controlled: $\nu = 0.52$, $\sigma_0 = 0.2$, $S_0 = 100$, $T = 1$, $r = 0$,
    $N = 252$, $M = 10{,}000$.
    Independent: $K$, $H$.
    Dependent: Asian call price.}
  \label{fig:sensitivitystrike}
\end{figure}

Figures~\ref{fig:sensitivity} and~\ref{fig:sensitivitystrike} show the
results:

\begin{itemize}
  \item \textbf{Price decreases with $H$, and more so at large $\nu$:} at
        $\nu = 0.52$ the ATM price falls from $5.373$ at $H = 0.05$ to $4.735$ at
        $H = 0.50$.  The gap between $H = 0.10$ and $H = 0.50$ is only $0.020$ at
        $\nu = 0.10$ but $0.937$ at $\nu = 0.70$: without vol-of-vol there is no
        volatility path whose roughness could matter.  As §\ref{sec:levy} shows,
        most of this effect is roughness rather than added variance.

  \item \textbf{Price increases with $\nu$:} at $H = 0.10$ the ATM price rises from
        $4.540$ at $\nu = 0.10$ to $5.869$ at $\nu = 0.70$.  Larger vol-of-vol
        amplifies $\sigma_t$ variation, and the increase is convex in $\nu$ because
        $\sigma_t = \sigma_0 e^{\nu W_t^H}$ enters exponentially.

  \item \textbf{The $H$ effect peaks at the money:} the spread between $H = 0.05$
        and $H = 0.50$ is $0.27$ at $K = 80$, $0.64$ at $K = 100$, and $0.41$ at
        $K = 120$ (Figure~\ref{fig:sensitivitystrike}), matching the strike profile
        of the premium in Figure~\ref{fig:levy}(b).
\end{itemize}

At the model parameters $H = 0.10$, $\nu = 0.52$ the ATM price in the grid is $5.223$,
one standard error ($\approx 0.095$) below the $10^6$-path reference of $5.325$
(§\ref{sec:convergence}); the grid's absolute level carries that single run's MC error,
while its differences do not.

\paragraph{Limitations: }
All grid points share one random seed.  This makes the trends across the grid
reliable, but the whole surface carries a common MC offset with a standard error of
$\approx 0.1$ price units.  The sensitivity surfaces are computed on a $6 \times 4$ grid
($H \times \nu$) at a fixed $N = 252$; finer grids and higher $M$ are feasible but
were not run for this report.
